% The abstract transfer has explicit construction hypotheses. Its final
% corollaries instantiate them with the families in the preceding sections.
\section{A finite transfer estimate}
\label{sec-finite-transfer}

We now estimate the length supplied by the assembly construction. The
starting size will be close to $n$, and most extensions will use the
families of Proposition~\ref{prop-global-counts}. A final sequence of
extensions of size three reduces the pending fraction far enough for
completion. We first give the transfer with explicit hypotheses and then
verify them using the constructions in the preceding sections.

\subsection{The construction statements needed for the estimate}

For a stage of old size $m\ge2$, retain the counts $R,Q,c,P$ from
Section~\ref{sec-assembly}. Define the rational numbers
\[
 y=\frac{Q}{m!},\qquad
 z=\frac{c}{(m-1)!},\qquad
 p=\frac{P}{(m-1)!},\qquad X=y+z.
\]
There are $(m-1)!$ cyclic classes on $m$ letters, so $p$ is the fraction
of classes still awaiting an assignment. Normalizing the other counts by
the indicated factorials will make inherited contributions unchanged
under an extension.

An extension from size $m$ to size $m+v$, where $v$ is a positive integer,
uses nonnegative integer counts $B,q,c_{\mathrm{adv}}$. Put
\[
 D_{m,v}=\prod_{i=0}^{v-1}(m+i),\qquad
 A_{m,v}=\prod_{i=1}^{v}(m+i).
\]
Lemma~\ref{lem-assembly-width} gives the bounds and update rules
\begin{align}
 0&\le B\le D_{m,v}, \label{eq-transfer-source-bound}\\
 R'&=A_{m,v}R+(m+v)PB,\notag\\
 Q'&=A_{m,v}Q+(m+v)Pq,\notag\\
 c'&=D_{m,v}c+Pc_{\mathrm{adv}},\notag\\
 P'&=P(D_{m,v}-B). \label{eq-transfer-integer-update}
\end{align}
The distinctness hypothesis in Lemma~\ref{lem-assembly-universe} is
required for \eqref{eq-transfer-source-bound} and for the complement
count in the last line.

Assume that the construction can start at any size $m_0\ge11$ with
\[
 (R,Q,c,P)=(0,0,0,(m_0-1)!).
\]
Use also the empty extensions from Section~\ref{sec-assembly}, which have
size one and $B=q=c_{\mathrm{adv}}=0$. At the final intermediate size
$h=n-3$, bounding the minimum in \eqref{eq-assembly-word} by $c+P$
gives an actual superpermutation $w$ with
\begin{equation}
 |w|\le F_3(n)+Q+(2h-1)(c+P).
 \label{eq-transfer-word-input}
\end{equation}
In the conditional results below, the initialization, the required
extensions, and completion to an actual superpermutation are all hypotheses. Equations
for counts without those construction statements are insufficient.

\begin{lemma}\label{lem-transfer-normalization}
The rules \eqref{eq-transfer-integer-update} preserve the identity
$R+mP=m!$. With
\[
 f=\frac{B}{D_{m,v}},\qquad
 g=\frac{q+c_{\mathrm{adv}}}{D_{m,v}},
\]
their normalized form is
\begin{equation}
 p'=p(1-f),\qquad
 y'=y+p\frac{q}{D_{m,v}},\qquad
 z'=z+p\frac{c_{\mathrm{adv}}}{D_{m,v}},\qquad
 X'=X+pg.
 \label{eq-transfer-normalized-update}
\end{equation}
In particular, $0\le p\le1$ throughout the construction, and empty
extensions preserve $y,z,p$.
\end{lemma}

\begin{proof}
Cancellation in the products gives
\[
 (m+v)D_{m,v}=mA_{m,v},\qquad
 (m+v)!=A_{m,v}m!,\qquad
 (m+v-1)!=D_{m,v}(m-1)!.
\]
Consequently,
\[
 R'+(m+v)P'
 =A_{m,v}R+(m+v)PD_{m,v}
 =A_{m,v}(R+mP).
\]
This is $(m+v)!$ when $R+mP=m!$. The initial counts satisfy that identity.
Dividing the updates for $Q',c',P'$ by the appropriate new factorials
proves \eqref{eq-transfer-normalized-update}. Initially $p=1$.
The inequality $0\le f\le1$ therefore proves $0\le p\le1$ by induction.
For an empty extension, $f=g=q=c_{\mathrm{adv}}=0$, so the normalized
quantities do not change.
\end{proof}

\subsection{Two finite inequalities}

The main estimate uses two short induction arguments. They apply to
rational numbers and do not require an approximation to a product.

\begin{lemma}\label{lem-transfer-pending}
Let $f_0,\ldots,f_{J-1}$ lie in $[0,1]$. Set $p_0=1$ and
$p_{i+1}=p_i(1-f_i)$. Then
\[
 p_J\left(1+\sum_{i=0}^{J-1}f_i\right)\le1.
\]
\end{lemma}

\begin{proof}
Write $s_j=\sum_{i<j}f_i$. All the numbers $p_j,s_j$ are nonnegative.
Using $s_{j+1}=s_j+f_j$, we obtain
\begin{align*}
 p_{j+1}(1+s_{j+1})
 &=p_j(1-f_j)(1+s_j+f_j)\\
 &=p_j(1+s_j)-p_jf_j(s_j+f_j)\\
 &\le p_j(1+s_j).
\end{align*}
The initial value is $p_0(1+s_0)=1$. Induction proves the claim.
\end{proof}

For example, if $f_0=f_1=1/2$, then $p_2=1/4$, and the left side of
the lemma is $1/2$. The estimate remains valid if a selected density is
one, in which case all subsequent pending mass is zero.

\begin{lemma}\label{lem-transfer-cost}
Suppose $p_{i+1}=p_i(1-f_i)$ and $X_{i+1}=X_i+p_i g_i$, with
$p_i\ge0$. If $t g_i\le Kf_i$ for every $i<J$, where $t>0$, then
\[
 t(X_J-X_0)\le K(p_0-p_J).
\]
\end{lemma}

\begin{proof}
For each $i$,
\[
 t(X_{i+1}-X_i)=tp_i g_i
 \le Kp_i f_i=K(p_i-p_{i+1}).
\]
Summing cancels all intermediate $X_i$ and $p_i$ and gives the result.
\end{proof}

\subsection{The schedule and its cost}

Fix an integer $r\ge6$ divisible by six. Let $W$ be a positive multiple
of $r$, let $E>0$, let $\rho\ge r^2$, and put
\[
 c_r=\frac{2(r-1)}{r-2}.
\]
At every multiple $m$ of $r$ with $m\ge4E$, assume that the construction
provides an extension of size $W$ for which
\begin{equation}
 f_m=\frac{\rho W}{m}(1+\theta_m),\qquad
 g_m=\frac{c_r\rho W}{m^2}(1+\varepsilon_m),\qquad
 |\theta_m|,|\varepsilon_m|\le\frac{E}{m},
 \label{eq-transfer-sparse-input}
\end{equation}
and
\begin{equation}
 c_{\mathrm{adv},m}\le\frac{r-1}{m}q_m.
 \label{eq-transfer-circle-input}
\end{equation}
The final extensions come from Proposition~\ref{prop-assembly-dense}.
For an extension of size three at any $a\ge49$, its inequalities and
Lemma~\ref{lem-transfer-normalization} give
\begin{equation}
 p'\le\frac34p,\qquad
 X'-X\le\frac8a(p-p'),\qquad
 z'-z\le\frac{8p}{a(a+1)}.
 \label{eq-transfer-dense-input}
\end{equation}
The extensions of size $W$ select a small proportion of the available
classes. The extensions of size three remove at least one quarter of
the pending mass. We use the latter only at the end.

Here is the derivation of \eqref{eq-transfer-dense-input}. Its selected
fraction is $f=B_a/D_{a,3}\ge1/4$, so
$p'=p(1-f)\le3p/4$. The cost estimate of
Proposition~\ref{prop-assembly-dense} gives
\[
 X'-X=p\frac{q_a+c_{\mathrm{adv},a}}{D_{a,3}}
 \le\frac8a\frac{pB_a}{D_{a,3}}
 =\frac8a(p-p').
\]
Finally, its circle count gives
\[
 z'-z=p\frac{(a+2)\nu_a}{a(a+1)(a+2)}
 =\frac{p\nu_a}{a(a+1)}\le\frac{8p}{a(a+1)}.
\]
These are updates of the existing stage. In particular, the final
extensions retain the rows and safe circles obtained during the sparse
part of the construction.

For a positive integer $n$, let $\ell$ be the least nonnegative integer
such that $n<2^\ell$. Define
\[
 K=12\ell,\qquad \tau=3+36\ell.
\]
For example, $\ell=7$ when $n=64$, because the upper inequality in the
definition is strict.

\begin{theorem}\label{thm-transfer-finite}
Assume the construction and completion statements above. If
\begin{equation}
 n\ge8rE,\qquad n\ge96r^2,\qquad
 r+W+\tau\le\frac{n}{2r},
 \label{eq-transfer-room}
\end{equation}
then there is a superpermutation $w$ on $n$ letters such that
\[
 |w|\le F_3(n)+\left(2+\frac{74}{r}\right)(n-4)!.
\]
\end{theorem}

\begin{proof}
Put $\alpha=(r-1)/r$, and let $m_0$ be the least multiple of $r$
that is at least $\alpha n$. Thus
\[
 \alpha n\le m_0<\alpha n+r.
\]
Set $a=n-\tau$ and
\[
 J=\left\lfloor\frac{a-m_0}{W}\right\rfloor,
 \qquad b=m_0+JW.
\]
Since $\lfloor x\rfloor>x-1$, the last condition in
\eqref{eq-transfer-room} gives
\[
 J>\frac{n/r-r-\tau-W}{W}
 \ge\frac{n}{2rW}>0.
\]
In particular $J\ge1$. Starting at $m_0$, use $J$ extensions of size
$W$, then $a-b$ empty extensions, and finally $K$ extensions of size
three. The floor defining $J$ gives $0\le a-b<W$. The final size is
$a+3K=n-3$, as required for completion.

The initial size satisfies $m_0\ge n/2\ge48r^2\ge11$. Each start
$m_i=m_0+iW$ of an extension of size $W$ is a multiple of
$r$ and lies between $n/2$ and $n$. Since $n\ge8rE$, it is at least
$4E$. Also $a\ge n/2$, because $\tau\le n/(2r)$, and
$a\ge n/2\ge48r^2>49$. Thus every extension used in this schedule
is within its stated range. The last extension of size three starts at
$n-6$.

Write
\[
 \delta=\frac{2E}{n}\le\frac1{4r}.
\]
The two relative errors at each $m_i$ have absolute value at most
$\delta$. In particular $f_{m_i}>0$. From
\eqref{eq-transfer-sparse-input},
\[
 \frac{ng_{m_i}}{f_{m_i}}
 \le\frac{c_r}{\alpha}\frac{1+\delta}{1-\delta}.
\]
Here
\[
 \frac{c_r}{\alpha}=\frac{2r}{r-2}
 =2+\frac4{r-2}\le2+\frac6r,
\]
because $r\ge6$. Since $0\le\delta\le1/2$,
\[
 \frac{1+\delta}{1-\delta}
 =1+\frac{2\delta}{1-\delta}
 \le1+4\delta\le1+\frac1r.
\]
Finally,
\[
 \left(2+\frac6r\right)\left(1+\frac1r\right)
 =2+\frac8r+\frac6{r^2}
 \le2+\frac9r.
\]
Let $p_s,X_s,y_s,z_s$ be the normalized values after the extensions of
size $W$. Lemma~\ref{lem-transfer-cost}, with $t=n$, now gives
\begin{equation}
 nX_s\le\left(2+\frac9r\right)(1-p_s)
 \le2+\frac9r.
 \label{eq-transfer-sparse-cost}
\end{equation}

We next bound $p_s$. Each selected density satisfies
\[
 f_{m_i}\ge\frac{\rho W}{n}(1-\delta)
 \ge\frac{\rho W}{2n}.
\]
The lower bound on $J$ therefore yields
\[
 \sum_{i<J} f_{m_i}
 \ge\frac{n}{2rW}\frac{\rho W}{2n}
 =\frac{\rho}{4r}\ge\frac r4.
\]
Lemma~\ref{lem-transfer-pending} proves
\begin{equation}
 p_s\le\frac4{r+4}\le\frac4r.
 \label{eq-transfer-pending-small}
\end{equation}
The intervening empty extensions preserve these normalized quantities.

All the final extensions start at sizes at least $a$. Their cost bound
in \eqref{eq-transfer-dense-input} telescopes to
\[
 X_{\mathrm{end}}-X_s
 \le\frac8a(p_s-p_{\mathrm{end}})
 \le\frac{8p_s}{a}.
\]
Using $n/a\le2$, \eqref{eq-transfer-sparse-cost}, and
\eqref{eq-transfer-pending-small}, we obtain
\begin{equation}
 nX_{\mathrm{end}}\le2+\frac{73}{r}.
 \label{eq-transfer-total-cost}
\end{equation}

The completion estimate also involves $z_{\mathrm{end}}$. During the
extensions of size $W$, \eqref{eq-transfer-circle-input} gives
\[
 z_s\le\frac{r-1}{m_0}y_s.
\]
Their local cost satisfies $m_i g_{m_i}\le8f_{m_i}$, since
$m_i\le n$ and $2+9/r\le8$. Applying
Lemma~\ref{lem-transfer-cost} with $t=m_0$ gives
$y_s\le X_s\le8/m_0$. Hence
\[
 z_s\le\frac{8(r-1)}{m_0^2}.
\]
For the final extensions, write $p_j$ for the pending mass before the
$j$th one. The contraction $p_{j+1}\le3p_j/4$ implies
$p_j\le4(p_j-p_{j+1})$. Summing gives
$\sum_{j<K}p_j\le4(p_s-p_{\mathrm{end}})\le4$.
The last inequality in \eqref{eq-transfer-dense-input} consequently gives
\[
 z_{\mathrm{end}}-z_s
 \le\frac{32}{a(a+1)}.
\]
Both $m_0$ and $a$ are at least $n/2$, so
\begin{equation}
 z_{\mathrm{end}}\le\frac{32r+96}{n^2}.
 \label{eq-transfer-final-circles}
\end{equation}

There are $12\ell$ final extensions. Since $p_s\le1$ and
$(3/4)^4=81/256<1/2$, their contraction gives
\begin{equation}
 p_{\mathrm{end}}
 \le(3/4)^{12\ell}<2^{-3\ell}<n^{-3}.
 \label{eq-transfer-final-pending}
\end{equation}
The strict inequalities are valid because $\ell\ge1$ and $n<2^\ell$.

Put $h=n-3$. Divide \eqref{eq-transfer-word-input} by $(n-4)!=(h-1)!$.
The normalized excess is at most
\begin{align*}
 h y_{\mathrm{end}}+(2h-1)(z_{\mathrm{end}}+p_{\mathrm{end}})
 &=hX_{\mathrm{end}}+(h-1)z_{\mathrm{end}}
   +(2h-1)p_{\mathrm{end}}\\
 &\le nX_{\mathrm{end}}+nz_{\mathrm{end}}+2np_{\mathrm{end}}.
\end{align*}
Equations \eqref{eq-transfer-total-cost}--\eqref{eq-transfer-final-pending}
bound this by
\[
 2+\frac{73}{r}+\frac{32r+96}{n}+\frac2{n^2}.
\]
For $r\ge6$, we have $32r+96\le48r$. The assumption $n\ge96r^2$
therefore bounds the penultimate term by $1/(2r)$. It also implies
$n^2\ge4r$, so the last term is at most $1/(2r)$. Their sum is at most
$1/r$, which proves the theorem.
\end{proof}

\subsection{One explicit threshold}

Choose the parameters in \eqref{eq-global-parameters}. By
Lemma~\ref{lem-polynomial-width}, they satisfy
$r^2\le W\le352r^9$ and $r^2\le\rho<r^2+1$. We need one further
consequence of those bounds.

\begin{lemma}\label{lem-transfer-parameter-bounds}
The integer $W$ is divisible by $r$, and
$E\le10^{12}352^4r^{58}$.
\end{lemma}

\begin{proof}
The identity $W=r^2A$ proves divisibility. The upper bound for $W$ gives
\[
 E=10^{12}r^{22}W^4
 \le10^{12}r^{22}(352r^9)^4
 =10^{12}352^4r^{58}.
\]
\end{proof}

\begin{lemma}\label{lem-transfer-binary-room}
For these parameters, the condition $n\ge8rE$ implies every condition
in \eqref{eq-transfer-room}.
\end{lemma}

\begin{proof}
Since $W\ge r^2\ge1$, we have $E\ge W$ and
$E\ge10^{12}r^{22}\ge2592r$. Consequently
\[
 n\ge8rE\ge(144r)^2.
\]
This proves $n\ge64$ and $n\ge96r^2$.

Recall that $\ell$ is the least integer with $n<2^\ell$.
Thus $\ell\ge7$ and $2^{\ell-1}\le n$. For every integer $t\ge7$,
\[
 t^2\le2^{t-1}.
\]
Indeed, the initial case is $49\le64$, and
$(t+1)^2\le2t^2$ for $t\ge7$, which proves the induction step.
It follows that $\ell^2\le n$. If $\ell\le144r$, then
$144r\ell\le(144r)^2\le n$. If $\ell\ge144r$, then
$144r\ell\le\ell^2\le n$. In either case,
\[
 36\ell\le\frac{n}{4r}.
\]
Finally, $W\ge r^2\ge r+3$, and therefore
\[
 r+W+3\le2W\le2E\le\frac{n}{4r}.
\]
Adding the last two inequalities gives
$r+W+\tau\le n/(2r)$, as required.
\end{proof}

\begin{corollary}\label{cor-transfer-polynomial-threshold}
For every integer $r\ge6$ divisible by six and every integer
$n\ge Cr^{61}$, there is an actual superpermutation $w$ satisfying
\[
 r|w|\le rF_3(n)+(2r+74)(n-4)!.
\]
\end{corollary}

\begin{proof}
Use the parameters of \eqref{eq-global-parameters}. At each multiple
$m$ of $r$ with $m\ge4E$, Proposition~\ref{prop-global-counts}
supplies actual source families above any chosen old cyclic order. Its
distinct source classes lie in the universe of
Lemma~\ref{lem-assembly-universe}, so $0\le B\le D_{m,W}$ and the
required complement exists. The same proposition supplies the literal
full run needed by Lemma~\ref{lem-assembly-connectors}. Together with
Lemma~\ref{lem-assembly-orbits}, this gives the safe supported circle
lists and transported components required by
Lemma~\ref{lem-assembly-width}.

Initialization and empty extensions are provided in
Section~\ref{sec-assembly}. The final extensions and their normalized
bounds are provided by Proposition~\ref{prop-assembly-dense} and
\eqref{eq-transfer-dense-input}. The schedule has final old size $n-3$,
and $n\ge Cr^{61}\ge11$, so \eqref{eq-assembly-word} gives its actual
superpermutation. This verifies all the construction and completion
hypotheses of Theorem~\ref{thm-transfer-finite}.

Lemma~\ref{lem-transfer-parameter-bounds} gives
\[
 8rE\le Cr^{59}\le Cr^{61}\le n.
\]
Lemma~\ref{lem-transfer-binary-room} supplies the other size conditions.
Apply Theorem~\ref{thm-transfer-finite} and multiply its conclusion by
the positive integer $r$.
\end{proof}

\subsection{The exponent in the asymptotic statement}

\begin{corollary}\label{cor-transfer-asymptotic}
\[
 L(n)=F_3(n)+2(n-4)!+
 O\!\left(\frac{(n-4)!}{n^{1/61}}\right).
\]
\end{corollary}

\begin{proof}
For $n\ge C6^{61}$, choose the largest positive multiple $r$ of six
with $Cr^{61}\le n$. Such an $r$ exists, and maximality gives
\[
 Cr^{61}\le n<C(r+6)^{61}\le C(2r)^{61}.
\]
The last inequality uses $r\ge6$. Taking positive sixty-first roots
now gives
\[
 \frac1r<2C^{1/61}n^{-1/61}.
\]
The word from Corollary~\ref{cor-transfer-polynomial-threshold} yields
\[
 \frac{L(n)-F_3(n)}{(n-4)!}-2
 \le\frac{74}{r}
 <148C^{1/61}n^{-1/61}.
\]
On the other side, use $(n-4)!=(n-4)(n-5)!$ to divide the imported lower bound
\eqref{eq-main-lower} by $(n-4)!$. Since $n\ge8$, this gives
\[
 2-\frac{L(n)-F_3(n)}{(n-4)!}
 \le\frac2{n-4}\le\frac4n\le4n^{-1/61}.
\]
Both sides are bounded by a constant times $n^{-1/61}$. Multiplying
by $(n-4)!$ proves the stated remainder.
\end{proof}
